\chapter{Sources}

\section{Speakers}

\textbf{These languages belong to the people who speak them, and none of this work exists without them.} The corpus everything else here is measured against is their words, written down.

The chapter listing the gold standard corpus opens every entry with the speaker and carries the conditions they set. A speaker's name is stated there once and nowhere else.

Most of the papers name their speakers. The rest cite a published dictionary and never say who spoke, and their entries say so. A linguist wrote the paper and a person transcribed the paper into a table, and neither of those is whose language it is.

\section{Papers with an automatic extractor}

Each of the papers below has an automatic extractor written for it and a hand extraction beside it. Every token of the language in each is accounted for in the corpus built from it, by the coverage check. For the two Lyon papers that statement is about the extracted text, the font's own encoding, and not what the page prints; the section below on those two says what that means.

\tiny
\setlength{\tabcolsep}{2pt}
\begin{longtable}{@{}>{\raggedright\arraybackslash}p{\dimexpr\textwidth*3/4-2\tabcolsep\relax}>{\raggedright\arraybackslash}p{\dimexpr\textwidth/4-2\tabcolsep\relax}@{}}
\toprule
Paper & Language \\
\midrule
Three Glossed Nɬeʔkepmxcín Narratives by Kʷəɬtəzétkʷu (Bernice Garcia) \cite{src:Garcia-Hannon-and-Stacey-2024} & nɬeʔkepmxcín \\
ɬ cutés us ɬ qəɬmín ɬ tmíxʷ (When Old One Created the Earth) \cite{src:Hall-and-Phillips-2025} & nɬeʔkepmxcín \\
Four Stories by wlwlmelst \cite{src:LaFontaine-and-Janzen-2024} & nɬeʔkepmxcín \\
Cw7aoz káti7 láti7 ku naxwít \cite{src:Matthewson-and-Redan-2026} & St'át'imcets \\
I Tsícwas sQwa7yán'ak Áku7 Graveyard Valley \cite{src:Alexander-and-Davis-2026} & St'át'imcets \\
Mary George Personal Narratives \cite{src:John-Hamilton-Davis-2021} & St'át'imcets \\
A Bella Coola tale: The Frog Children \cite{src:Nater-2015} & Nuxalk \\
Three Okanagan stories about priests \cite{src:Lyon-2015} & Nsyilxcən \\
12 more Upper Nicola Okanagan narratives \cite{src:Lindley-and-Lyon-2013} & Nsyilxcən \\
Poking Fun in Lushootseed \cite{src:Hilbert-1983} & Lushootseed \\
A Comparative Analysis of Stress in Northern and Southern Lushootseed \cite{src:Mellesmoen-and-Kye-2026} & Lushootseed \\
BC Indigenous people's Chinuk pipa script: History, analysis, and texts \cite{src:Robertson-2012} & nɬeʔkepmxcín, Secwepemctsín \\
Lexical Suffixes and Connectives in Proto-Central Salish and Beyond \cite{src:Wolfe-2025} & eighteen languages \\
Voiceless Words in Bella Coola: Fact vs. Fiction \cite{src:Nater-2024} & Nuxalk \\
Nsyilxcn Inchoatives and their Distributions Across Root Types \cite{src:Lyon-2025} & Nsyilxcən \\
The truncated reduplication in Twana \cite{src:Kim-2017} & Twana \\
How Salish is Bella Coola? \cite{src:Nater-2013} & Nuxalk \\
Ctrl-Alt-Delete: The Control Directive and Associated /t/ Deletion in nɬeʔkepmxcín \cite{src:Hall-Luntzlara-Mellesmoen-and-Reid-2026} & nɬeʔkepmxcín \\
\bottomrule
\end{longtable}
\normalsize

Each automatic extractor writes the marked record, the gold standard lines holding only target-language speech, a list of what it could not classify, and for the two Lyon papers the word forms recovered from the interlinear. A record is named speaker first, because the speaker is who the corpus is of.

The papers differ in what makes them hard to extract, and the differences are the evidence for how each extractor is built.

\subsection{Surface and underlying forms}

\textbf{Hall, Luntzlara, Mellesmoen and Reid print every form twice, and only one of the two was said} \cite{src:Hall-Luntzlara-Mellesmoen-and-Reid-2026}. They argue about what the nɬeʔkepmxcín transitive suffixes are underlyingly. Each example carries a surface form in square brackets and the underlying form they propose for it in slashes: \texttt{[wéc̓entis]} beside \texttt{/wéc̓e-n-t-ey-es/}. The hand extraction is 494 rows and verifies in both directions with nothing left over.

That layout makes the delimiter the evidence, and the extractor reads only the delimiter. A bracketed run is a word and reaches the gold standard corpus. A slashed run is the authors' analysis, holds morpheme boundaries and a null-morpheme sign, and does not. A star marks a form the analysis predicts and the language does not have, written \texttt{*[kícetxʷ]} in one place and \texttt{[*kícexʷ]} in another, and the extractor tests for both shapes. The delimiters stay on the form in the record, because the record holds what the page printed, and come off on the way to the gold standard corpus, where the entry is the word.

The extractor does not find two kinds of form, by design. Examples (24) and (29) to (31) are ordered-rule derivations and every line between the underlying form and the surface form is a stage: \texttt{kícntene}, \texttt{kícntne}, \texttt{kícnne}, \texttt{kícne}. Those carry no delimiter, because the paper sets them in a column. The prose also prints some surface forms bare, \texttt{kicnəxʷ} and \texttt{wíktxʷ} among them. The hand extraction has all of them and the extractor leaves them unclassified, which for a form nobody said is the safe direction to be wrong in.

Section 3.1.1 argues from two other languages, and the speaker column carries them: ʔayʔaǰuθəm has \texttt{-θi} and St'át'imcets has \texttt{-ci}, with Newman's proto-Salish \texttt{*c} and \texttt{*ci} cited beside them. The extractor puts zero forms under the wrong language. One string is genuinely ambiguous: \texttt{-ci} is both the St'át'imcets suffix the paper cites and the nɬeʔkepmxcín suffix the paper proposes, and nothing in the token separates them. \texttt{-ci} enters as nɬeʔkepmxcín, which the paper argues it is.

The extractor removes matched spans by position and not by name. Page 13 carries \texttt{/ʔ/}, which yields the one-character form \texttt{ʔ}, and removing that by name would also remove it from \texttt{\{n, n̓, h, ʔ\}]morpheme-t-\_\_\_}. One token does not come back through the coverage check: \texttt{ʔ\}]morpheme-t-\_\_\_} holds a brace, the record delimits its spans with braces, and the check splits it on the way back out. The token is in the paper and in the record, and the paper stands at 99.7 percent coverage for it.

\subsection{An etymological lexicon}

\textbf{Nater's \emph{How Salish is Bella Coola?} is the largest table in the set} \cite{src:Nater-2013}. Nater asks how Salish Bella Coola is, and his answer is his whole lexicon sorted by where each word came from: 209 proto-Salish, 94 Coastal, 63 Interior, 214 non-Salish, 34 areal, and 661 of unknown origin. That is 1407 numbered entries over 52 pages, more of one language than the rest of this set holds together.

The fourth column of its two tables is whichever language Nater is comparing against, and the extractor finds where that column starts by the token that opens it: one of the sixteen abbreviations section 1.1 defines, a language the body defines, or a reconstruction opening with an asterisk. Haisla, Heiltsuk, Oowekyala and Kwak̓wala are North Wakashan. Proto-Athabascan, Eyak, Carrier and Tahltan are Na-Dene. Nootka, Quileute, Chinook and Yurok are four families past that. Squamish, Sechelt, Shuswap, Lillooet and Halkomelem are Salish and still not this language. The speaker column carries all of it, and only Bella Coola reaches the gold standard corpus.

The paper prints that column's own language at the head of every table page, as \texttt{Line Gloss BC CS Cognates}, and the extractor reads it from there. A bare \texttt{*ƛ‟ǝp} under that head is Coast Salish and holds nothing in itself that an explicit tag would have supplied. Reading the head places 326 reconstructions under the name of the group they reconstruct.

\textbf{Section 4 prints each of its entries twice.} \texttt{AKW'A} and \texttt{ʔak‟ʷa} are one word, and section 2.1 gives the mapping: \texttt{ʷ} is written \texttt{w}, \texttt{c} is \texttt{ts}, \texttt{ƛ‟} is \texttt{tl‟}, \texttt{ɬ} is \texttt{lh}, \texttt{x} is \texttt{c}, \texttt{χ} is \texttt{x}, and \texttt{ʔ} is \texttt{7}. Both are Bella Coola and the table records both. Only the phonemic column reaches the gold standard corpus, because two orthographies in one stream would put the byte pairs of a transliteration into the measurement of a language.

\textbf{One paper can write one sound two ways.} The tables set every ejective as \texttt{‟}, U+201F, 1763 times. Section 3's prose sets it as \texttt{’}, U+2019, 51 times. The prose's \texttt{t’nχʷ} and the appendix's \texttt{t‟nχʷ} are one word in two encodings. The schwa does the same, \texttt{ǝ} U+01DD 183 times against \texttt{ə} U+0259 52 times, and NFC unifies neither pair. A mark set built from the tables alone reads section 3 as English.

Its practical orthography also writes the glottal stop as the digit \texttt{7}, which the shared mark set holds. A citation like \texttt{Sq67:137} or \texttt{Ku02:86-87} therefore carries a character of the language, and the second direction of the check asks about it. A digit that is a letter in one orthography cannot also be punctuation in a check that covers that orthography.

The table records each cognate as its own row, and the extractor writes the cognate column as the single string the paper printed: \texttt{*ƛ‟ǝp „deep (water)‟} is one item against two rows. The extractor's large counts of forms not found and forms added are that one structural difference counted from two sides.

\subsection{Twana}

Kim's paper on Twana reduplication \cite{src:Kim-2017} is the only Twana source here. Its hand extraction is 358 rows. Every Twana form in it is Drachman's, out of a 1969 dissertation the paper calls the only reliable reference in existence for Twana CVC reduplication. What is on these pages is close to everything written down.

Kim also carries the trap this set makes easiest to fall into. Example (1) sits where the data usually starts and it is Tillamook, in Edel's 1939 transcription with its own vowel letters; (19) and (20) are Thompson and Lillooet; footnote 9 mixes Puget Sound Salish, Moses-Columbian and a row of English phonetics. The speaker column carries all of it and the footnote is left unclassified, because the language changes inside one sentence there.

\textbf{The lateral fricative has a third character.} Papers differ between \texttt{ɬ} and \texttt{ł}, and the text space carries both. Kim writes it \texttt{ɫ}, U+026B, sixty-two times, and a mark set holding the other two reads every one of those tokens as English. The same paper writes the glottal stop two ways, \texttt{ʔ} where it is phonemic and \texttt{ˀ} where a rule derived it, which its footnote 11 states outright.

The extraction doubles a combining mark on four phonetic forms. Those four are corrected at the source, and the table records the form the page prints; transcribing the doubled form would put a corrupt string in the table.

\subsection{Elicited roots}

Lyon's survey of how 48 Nsyilxcn roots inchoativize \cite{src:Lyon-2025} is elicited from ɬk̓mxnalqs Delphine Derrickson-Armstrong and c̓əskʕáknaʔ Dave Michele of stq̓aʔtkʷɬniw̓t. Its hand extraction is 537 rows and verifies in both directions. Most cells of its two tables are starred, a form the linguist built and the speakers rejected, and those are held out of the gold standard corpus: a smaller corpus is better than one seeded with words the language does not have. The extractor files the Spokane, Secwepemctsín, Lillooet, nxaʔamxčín, Thompson and Coeur d'Alene cognates under their own languages.

The paper stands at 99.5 percent coverage, and neither of the two missing tokens is a word. Closing the space after a glottalization mark welds \texttt{√pul̕} and \texttt{púl̕} on page 9 onto the \texttt{qt:\{} that opens the dictionary quotation after them. The welded tokens hold a brace, the record delimits its spans with braces, and the coverage check splits them on the way back out. They exist in neither the page nor any corpus.

\subsection{Words with no vowel}

Nater's \emph{Voiceless Words in Bella Coola} \cite{src:Nater-2024} argues that Bella Coola has words with no vowel in them, and its data is 127 numbered entries that are each a run of two to four consonants. Six more entries and two comparison tables are Heiltsuk, Oowekyala, Kwak̓wala and Haisla, which are North Wakashan and not Salish at all, and the speaker column carries the distinction. Its hand extraction is 671 rows and verifies in both directions with nothing left over.

\subsection{A comparative reconstruction}

Wolfe's paper \cite{src:Wolfe-2025} is a comparative reconstruction. Instead of one speaker telling one story it cites lexical suffixes from published dictionaries of eighteen languages, laid out in columns under a two-letter abbreviation. Its hand extraction is 772 rows and verifies in both directions with nothing left over. Its wrong-language count carries the most weight of any paper here, because it holds more than two languages, and it reads zero: the speaker column carries the language on every row and the extractor keys its output to the same column.

\section{Hand extraction}

The automatic extractor is not the source of the corpus. Each paper is transcribed by a person into a table that says, for every form in the paper, what it is, and the extractor is graded against that table. The table pairs each form with its meaning; it is a glossed word list, and it is the research's own data.

Every table is checked against its paper in both directions: no form written that the paper does not hold, and no language token in the paper that no row covers. Of the 25 tables, 14 verify with nothing left over in either direction. Of the rest:

\begin{itemize}
\item two are checked against a drafted page text, because their extracted text is not the page, and what their disagreements measure is the damage in that text (the next section);
\item one, Robertson \cite{src:Robertson-2012}, is damaged four ways in its extraction, and its disagreements are the same finding counted from both sides;
\item five carry none of their own orthography in their text layer, the scanned 1975 typescript and the four ICSNL 2 papers, and are not counted;
\item three carry the failures the chapter \emph{How wrong it could be} counts: Givens and Hall \cite{src:Givens-and-Hall-2023} with 80 forms the paper does not hold, Nater's lexicon \cite{src:Nater-2013} with 81, and Alexander and Davis \cite{src:Alexander-and-Davis-2026} with 2.
\end{itemize}

How each automatic extractor fares against its table is reported per paper in the chapter \emph{How wrong it could be}, the only form in which those counts mean anything.

\textbf{One hole in the second direction is not a hole in a table.} Nater's argument is that Bella Coola has words with no vowel in them. Entries like \texttt{kp} `each, all, every', \texttt{tp} `spotted', \texttt{ps}, \texttt{px} and \texttt{sx} are plain ASCII with nothing in them that any mark set can catch. The language-token test needs a marked character, and about thirty of the 133 entries are invisible to the second direction. The first direction still checks every one of them, because it looks for what the table wrote down. The apostrophe cannot be the way in: Nater writes every ejective with one, but a gloss ends in a closing quote and the two characters are the same, and every English gloss then reads as a word of the language. That trade is a few hundred false holes for about thirty real ones, and the apostrophe stays out.

\textbf{Verification separates the author from the speaker.} In Hilbert's paper \cite{src:Hilbert-1983} the coverage check is 100 percent, and the hand extraction still corrects two things. Vi Hilbert wrote the paper; the twenty-one examples were said by her aunt \textbf{Susie Sampson Peter} of the Upper Skagit and by \textbf{Martha LaMont}, recorded by Leon Metcalf between 1950 and 1958 and by Thom Hess in 1963. And example 2's English translation is a single line, ``High class, high class was Raven.'' The extraction welds the next eleven lines of Hilbert's commentary onto it. The typescript sets examples in an indented column and the essay at full width; the extraction drops the indent and the width recovers it.

\textbf{Six of the PDFs break words in half.} They leave a space after a stacked diacritic: \texttt{K̓weswapáw̓} arrives as \texttt{K̓} and \texttt{weswapáw̓}. The counts are Hall and Phillips 996, Garcia 943, LaFontaine and Janzen 478, Mellesmoen and Kye 298, Matthewson and Redan 169, Mary George 159. The coverage check cannot see it, because it puts both sides through the same repair and a word broken on both sides matches itself. A repair closes it; \texttt{ʷ} is held out, because it is a spacing letter and a space after one is a real boundary.

The stress accents are held out of that repair. A word can end in a stressed vowel, and \texttt{ntes neʔé e sqyéytn} in LaFontaine and Janzen would close to \texttt{neʔée}, which the language does not have. Mellesmoen and Kye passes its own wider set, because that PDF prints two spaces at a real boundary and a lone space after an acute there is always the inserted one.

Matthewson and Redan carries a second form of the same damage that no rule can close safely. The PDF also breaks \emph{before} a marked letter: \texttt{cácl̓ep} arrives as \texttt{các l̓ep}. The same shape is a real word boundary at \texttt{lta q̓íl̓qa}, and nothing in the characters separates the two. The extractor leaves both alone; the hand extraction has the true forms.

\textbf{The apostrophe is a letter.} \texttt{’} is glottalization in Nuxalk and in Lyon's Okanagan, and stripping it as punctuation turns the enclitic \texttt{˽c’} into \texttt{˽c}.

\section{Papers whose extracted text is not the page}

Six of the 146 PDFs held carry a font that renumbers its glyph codes with an \texttt{/Encoding /Differences} array and declares no \texttt{/ToUnicode} map. An extractor is then handed code numbers with nothing to turn them into, reads them in a default encoding, and what can land in the text is the font's private alphabet.

Each of the six is opened against its own rendered pages, and they come out four different ways.

\tiny
\setlength{\tabcolsep}{2pt}
\begin{longtable}{@{}>{\raggedright\arraybackslash}p{\dimexpr\textwidth/3-2\tabcolsep\relax}>{\raggedright\arraybackslash}p{\dimexpr\textwidth/3-2\tabcolsep\relax}>{\raggedright\arraybackslash}p{\dimexpr\textwidth/3-2\tabcolsep\relax}@{}}
\toprule
Paper & Fonts declaring no map back to Unicode & What the page shows \\
\midrule
Lyon \cite{src:Lyon-2015} & NimbusRomNo9L Regu, Medi, ReguItal & damaged, and transcribed from the pages in full \\
Lindley and Lyon \cite{src:Lindley-and-Lyon-2013} & NimbusRomNo9L Regu, Medi, ReguItal & damaged the same way, and transcribed from the pages in full \\
Lyon \cite{src:Lyon-2014} & NimbusRomNo9L Regu, Medi, MediItal, ReguItal & damaged, by a different table \\
Robertson \cite{src:Robertson-2012} & Symbol and five TrueType subsets & damaged four ways, one of them decodable exactly \\
Abraham \cite{src:Abraham-2015} & NimbusRomNo9L Regu, Medi & clean \\
Lonsdale and Matsushita \cite{src:Lonsdale-and-Matsushita-2011} & Courier, Times-Roman, Times-Italic, CMSY10 & clean \\
\bottomrule
\end{longtable}
\normalsize

A missing \texttt{/ToUnicode} is not the fault by itself. 141 of the 146 hold a font without one and nearly all of them extract correctly, because a standard encoding already says what the codes mean and every extractor has that table. It is the renumbering that does the damage, and 6 of those 141 renumber.

Renumbering does not always damage the text. Two of the six extract correctly, and both are checked line for line against a rendered page. Abraham's paper is St'át'imc in the van Eijk orthography, which is ASCII apart from the accented vowels, and its renumbered codes never reach a character they would damage: page 1's first line reads \texttt{Ats’xenlhkán ta sásqets áku7 Nséq’a (Charlie Mack’s), lti Líl’wata Tsel’álh c.walh,} on the page and in the text. Lonsdale and Matsushita print their Lushootseed in the ASCII transliteration their own parser reads, and there is nothing on the page for a font to lose. Page 5 sets \texttt{LEFT-WALL ?u+ da?a +d ?ElgWE? ?E kWi s+ gWistalb ti?E? SukWE? .} and the extraction gives that string exactly.

Page 23 of Lindley and Lyon prints

\texttt{cítxʷsəlx uɬ t̓i nyʕ̓ip ck̓aʔítət}

and the extracted text holds

\texttt{cítxws@lx uì ’ti ny ’Qip c ’kaPít@t}

Every schwa is \texttt{@}, every lateral fricative \texttt{ì}, every pharyngeal \texttt{Q}, every glottal stop \texttt{P}; the wedge stands in front of its letter as \texttt{ˇx}, the ejective mark stands in front of its letter instead of over it, and the word \texttt{ck̓aʔítət} is split. Two independent PDF libraries lose the same things. This is then a property of the file and not of one library.

\textbf{One part of the loss cannot be inverted.} Labialization is written with a raised w and the page also has a plain w. Both arrive as \texttt{w}. Page \texttt{kʷukʷ} and page \texttt{wist} are the same string in the text, and nothing in the file separates them.

The consequence for method is the reason this section exists. A hand extraction taken from one of these texts records the font's encoding and not the paper, and it then verifies clean against that same text in both directions, because both sides of the check are reading the one damaged artifact. The verification of these papers is therefore done against the page.

The pages are rendered as images and the transcription is done from the image. A page text is drafted from the font's codes in the shared orthography, line for line with the extraction, and a page of one is then a page of the other. The draft is a starting point a person corrects against the rendered page, never a source. There are seven things it cannot decide:

\begin{itemize}
\item \textbf{Labialization}, for the reason above. The drafter writes \texttt{ʷ} after the consonants that take one, which is right for \texttt{kʷ qʷ xʷ x̌ʷ} and wrong wherever a real \texttt{w} follows a consonant. The rule errs in both directions, and page 315 has one of each. \texttt{púlxwi} gets a \texttt{ʷ} it should not have, and section 3.2 settles that by parsing the word \texttt{√púlx-wi}. \texttt{lkʷu·········t} loses the one it should have, because the inserted space splits it into \texttt{lk} and \texttt{wu·········t} before the rule can see the consonant, and section 3.2 parses that one \texttt{√lkʷ=ut}. The rule runs over the whole line and reaches the English as well: stanza 117's word gloss \texttt{he.fell.off.backwards} comes out \texttt{he.fell.off.backʷards}.
\item \textbf{Word boundaries.} The PDF inserts a space in front of a letter carrying a mark. \texttt{s ’plá ’ks@lx} is one word and \texttt{iP ’kl} is two, and both are a space in front of a marked letter. Page 25 settles the first as \texttt{sp̓lák̓səlx}, page 24 the second as \texttt{iʔ k̓l}.
\item \textbf{The space that was not inserted.} Where the PDF leaves the mark against its letter there is nothing to key on, and a medial \texttt{’} is also the apostrophe of \texttt{Lyon’s}. \texttt{nkw’ritkw}, \texttt{xw’tilx} and \texttt{wa’y} keep the mark standing in front of its letter; the pages read them \texttt{nkʷr̓itkʷ}, \texttt{xʷt̓ilx} and \texttt{way̓}.
\item \textbf{A P-initial word with no other mark on it.} \texttt{Pitx}, \texttt{Pums}, \texttt{Pistkm} and \texttt{Paksyilx} come through as English, because a capital P opening an otherwise unmarked token is the shape of \texttt{Papers} and \texttt{Penticton}.
\item \textbf{The dropped space.} At a change of font and at some line ends the PDF loses the space entirely. \texttt{Okanagan yaxʷt}, \texttt{uɬ ny}, \texttt{məɬ ixíʔ} and \texttt{axáʔ ikɬcítxʷ} arrive welded into one token.
\item \textbf{A run of the length mark before a closing quote.} The raised dot arrives as \texttt{;}, and a rule turns it back wherever a letter or a spacing mark follows it, because a lone \texttt{;} is also Lyon's semicolon. A run of eight is never punctuation, and \texttt{staʔx̌íləm “whoa········”} is the only place where a run is followed by a closing quote instead. That pattern occurs twice, both in the priests paper.
\item \textbf{A one-letter gloss label that is also a glyph code.} A gloss token is recognized by a run of two capitals, because one capital on its own is also \texttt{iP} and \texttt{Philosophical}. Stanza 108's gloss line labels the question marker \texttt{Q}, one capital, and the token comes out \texttt{ʕ}. Nothing inside the token separates the two readings. That label occurs twice, both in the priests paper.
\end{itemize}

Every disagreement the check reports on Lindley and Lyon falls in these classes, and the last two occur in the priests paper alone. That count is the measurement the draft exists to produce.

The extraction is better evidence than the rendered page for one thing. The transposition moves a mark off its letter and preserves it, and a token's mark count survives the channel exactly. At the scale these pages render, a run like \texttt{m̓y̓m̓y̓á} cannot be told from \texttt{m̓ym̓á} by eye, and reading the page alone undercounts. Stanza 112 arrives as \texttt{s- m̓ y̓- m̓ y̓-á y̓-s}, which is five marks, and the word is \texttt{s-m̓y̓-m̓y̓-áy̓-s}. The count is taken from the text and the boundaries from the page.

\subsection{A scanned typescript is a worse case than any font}

Hilbert and Hess \cite{src:Hilbert-and-Hess-1975} carries no font problem at all. It is a 1975 typescript, scanned, and its text is OCR of the scan. Page 1 prints

\texttt{kʷa? ləcux̌əčadadid dxʷ?al ti?ə? s?acuss (h)əlgʷə? dəxʷx̌əbil ?ə ti?ə?.}

and the text holds

\texttt{kWa? lacuXacadadid dxW?al ti?a? s?acuss (h)algWa? daxwxabil ?a ti?a?}

Every schwa is \texttt{a}, every raised \texttt{ʷ} is \texttt{W}, the wedge over \texttt{x} is gone, and \texttt{sudᶻigʷicuts} comes through as \texttt{sudZigWiOltS}.

The difference from the Lyon papers matters more than the size of it. There the damage is a font's renumbering, which is deterministic: one code always stands for one glyph, which lets a table invert it. Here the damage is OCR guessing at shapes on a scan, and \texttt{sudᶻigʷicuts} to \texttt{sudZigWiOltS} is not a substitution anything can run backwards. The hand extraction is not a check on the extraction for this paper; it is the only correct record of it.

That also puts a floor under the whole method. A paper this old cannot be verified in both directions, because one direction has no sound source to verify against. What the check reports on it is the distance between a page and its OCR, which is worth having and is not the same measurement it makes elsewhere.

The check strips \texttt{?} as punctuation, and this orthography and the 1983 typescript's both write the glottal stop as \texttt{?}. \texttt{?ə} bares to \texttt{ə}, and a word with a glottal stop cannot be told from one without. The answer has to be per paper, since \texttt{?} is punctuation in every other paper here, and it is not applied.

Page 6 of that paper is a comparative appendix in Upper Chehalis and Cowlitz, thirteen sentences of two other Salishan languages. The table names each of them in its own column, which keeps them out of a Lushootseed corpus.

\subsection{The damage is a mutation distribution}

The damage is a hazard to route around, and it is also the one calibrated mutation sample in the archive. The damaged text is kept for that reason.

Two encodings of one text, aligned line for line, with the page settling which is right, is a measured channel. The extracted text is the mutated side, the drafted page text the other, and the hand extraction from the rendered page is ground truth for both. No other paper here has a second encoding of the same page to compare against.

The mutations sort into five kinds, and the last two are what make this a distribution instead of a rewrite table.

\tiny
\setlength{\tabcolsep}{2pt}
\begin{longtable}{@{}>{\raggedright\arraybackslash}p{\dimexpr\textwidth/3-2\tabcolsep\relax}>{\raggedright\arraybackslash}p{\dimexpr\textwidth/3-2\tabcolsep\relax}>{\raggedright\arraybackslash}p{\dimexpr\textwidth/3-2\tabcolsep\relax}@{}}
\toprule
Kind & Instances & Invertible \\
\midrule
Substitution, one symbol for one & \texttt{ə}→\texttt{@}, \texttt{ɬ}→\texttt{ì}, \texttt{ʕ}→\texttt{Q}, \texttt{ʔ}→\texttt{P}, \texttt{ƛ}→\texttt{ň}, \texttt{·}→\texttt{;} & yes \\
Collapse, two symbols onto one & \texttt{ʷ}→\texttt{w}, and plain \texttt{w}→\texttt{w} & no \\
Transposition & a combining mark over its letter becomes a spacing mark before it: \texttt{k̓}→\texttt{’k}, \texttt{x̌}→\texttt{ˇx} & yes \\
Insertion & a space appears in front of a letter carrying a mark & not decidable alone \\
Deletion & the space at a change of font is dropped, welding two words & not decidable alone \\
\bottomrule
\end{longtable}
\normalsize

The insertion does not always fire. Lindley and Lyon's text holds \texttt{wa ’y} in most places and \texttt{wa’y} in three, for the single word the page prints as \texttt{way̓}. Its rate lies strictly between 0 and 1, and the same input shape has two outputs. The insertion fires with a probability, and a repair applied without the page welds words that were never one word.

The deletion fires where the paper changes font mid-sentence. Footnote 5 of the same paper sets its cited forms in italic inside roman prose, and the two arrive as \texttt{Okanaganyaxʷt} and \texttt{Okanaganyaʕyáʕt}, each one token where the page prints two words. The second is recoverable, because \texttt{yaʕyáʕt} is a common word and appears on its own elsewhere in the paper. The first is not: \texttt{yaxʷt} occurs nowhere else. Which of the two happens depends on the rest of the corpus and not on the token, the second reason this is a distribution.

The collapse sets a ceiling. No recovery separates page \texttt{kʷukʷ} from page \texttt{wist} in the mutated text, because the channel destroys the distinction instead of disguising it. Anything downstream that claims to have recovered a labialized consonant on these two papers is claiming something the file cannot support.

A detector keyed to a glyph inventory reads the mutated side as holding no language at all, and reports a paper with eleven unread texts as fully covered. A detector reading the distribution of symbols survives the same input. The gold standard corpus says how well a vector separates the target language from the English around it; this pair says how much mutation it takes before it stops separating them at all.

The per-symbol rates are not measured. A candidate mapping is scored by how many mutated tokens become forms attested in a clean paper, which is a proxy for the answer. The page-verified table is the answer itself for the lines it covers.

\textbf{The conversion table is per document and not per font family.} Four of the six papers above are set in NimbusRomNo9L, and the two with page-verified tables supply a candidate mapping for the other two. Neither takes it, for two different reasons. Abraham's paper has nothing to convert. Lyon's paper on the linguistic evidence for a Francis Drake landing in Oregon \cite{src:Lyon-2014} is damaged and wants a different table: its language data comes from Frachtenberg's Oregon transcriptions, and while \texttt{ì} for \texttt{ɬ} carries over and so does the transposed acute, it also needs a transposed macron the Okanagan papers never use, \texttt{E} for \texttt{ɛ}, and several more marks. One code disagrees outright. \texttt{;} is the raised length dot in the Okanagan papers and the glottalization mark here. Page 24's \texttt{kʼeuʼts!} arrives as \texttt{k;eu´ts!}, and the Okanagan table would silently turn every glottalization into a length mark. Scored, it puts 4 of 407 occurrences into attested forms against 5 before it.

Robertson's paper \cite{src:Robertson-2012} is damaged four ways, and one of them costs nothing. Its extractor could not resolve a glyph and printed the glyph's name instead. The text holds \texttt{/uni0294oo /uni026C /xé/uni0294} where page 30 sets a morphemic line. A \texttt{/uniXXXX} name carries the code point it stands for: \texttt{/uni0294} is \texttt{ʔ}, \texttt{/uni026C} is \texttt{ɬ}, \texttt{/uni0259} is \texttt{ə} and \texttt{/uni019B} is \texttt{ƛ}. All 103 of them are turned back, along with \texttt{/combiningdotbelow} and \texttt{/combiningacuteaccent}, and no glyph name is left standing anywhere in the paper. That part is arithmetic and exact. The other three come off the pages:

\begin{itemize}
\item \textbf{Labialization, collapsed exactly as it is in the Lyon papers.} Page 30 sets \texttt{/kʷú[·kʷ]piʔ} and the text gives \texttt{/kwú[·kw]pi/uni0294}. The paper's own prose loses it too: page 30 names its symbols \texttt{(č, š, xʷ)} and the text holds \texttt{( č, š, x w)}, with a space where the raised letter was.
\item \textbf{A word broken across two lines with no hyphen.} 42 lines of the extraction hold nothing but a fragment: \texttt{Da} above \texttt{vid D. Robertson}, \texttt{ma} above \texttt{ximal efficiency}. Most of them are in the interlinear, where the break falls inside a gloss: \texttt{A} above \texttt{UG-}, \texttt{s} above \texttt{ick}, \texttt{gr} above \texttt{eet}. A rule that joins every fragment line to the one under it would be wrong in the alphabet tables, where \texttt{wa}, \texttt{wi} and \texttt{waw} are Chinuk pipa letter names and each is a line of its own.
\item \textbf{A glyph repeated where the page prints it once.} Page 1's epigraph prints \texttt{t’íxʷəɬ γ-ʔ-[χ-]q’y-n’-tén} with one \texttt{ɬ} and one \texttt{ʔ}, and the text holds four of each. The doubling happens on 2 lines in the paper and both are set in bold italic.
\end{itemize}

The hand extraction is 195 rows and covers both Salish texts whole, the Chinuk pipa and Chinook Jargon forms cited through the analysis, and the two bibliography entries written in Chinuk Wawa. Texts 3 through 6 are Chinook Jargon, which is a pidgin; they are the paper's own material and they are not Salish, and the speaker column keeps them out of the target corpus.

The extractor finds all 20 stanzas of Text 1 and all 7 of Text 2 with their numbering intact. The transliteration row is plain ASCII: \texttt{o l ha l kukpi,} carries no character of the modern orthography, and the test every other extractor leans on would call it English. The row is identified by where it sits. The gold standard corpus takes the transliteration row alone: the morphemic row is Robertson's analysis and holds forms nobody said. The line join is applied identically by the extractor and the coverage check, and with it the paper is at 100 percent coverage.

\subsection{The per-run test reads the damage a blanket repair cannot}

Reading a broken word off a rendered page works and does not scale: twenty papers at twenty-four pages each is four hundred images, and every reading is a judgment nobody can check afterward. The screening can be asked instead.

A break site is a place where the extraction has two tokens and the paper may have had one. The candidates are the two joined by each combining mark the paper writes, and joined by nothing, the reading where the space is a real boundary. A candidate counts as attested where the paper writes it as a token, and where the paper writes it inside one: \texttt{qʷal út} is \texttt{qʷal̓út}, which Davis and Mellesmoen \cite{src:Davis-and-Mellesmoen-2023} never print alone and do print inside \texttt{qʷəqʷal̓út}. Scoring is the per-run test of the dialect border, with run counts taken from the tokens carrying no break and flattened to maximum entropy. A run common everywhere contributes nothing, and what is left belongs to that restoration.

\textbf{It finds 86 sites across the fifteen papers with extractors, and every one of those papers reports 100 percent coverage.} That is not a contradiction: the coverage check puts both sides through the same repair, and a word broken in the source and broken in the extraction matches itself.

\tiny
\setlength{\tabcolsep}{2pt}
\begin{longtable}{@{}>{\raggedright\arraybackslash}p{\dimexpr\textwidth/3-2\tabcolsep\relax}>{\raggedright\arraybackslash}p{\dimexpr\textwidth/3-2\tabcolsep\relax}>{\raggedright\arraybackslash}p{\dimexpr\textwidth/3-2\tabcolsep\relax}@{}}
\toprule
paper & read & declined on score \\
\midrule
Garcia, Hannon and Stacey & 24 & 9 \\
Hall and Phillips & 24 & 0 \\
LaFontaine and Janzen & 18 & 2 \\
Davis and Mellesmoen & 4 & 1 \\
Kim & 4 & 0 \\
Mary George & 4 & 0 \\
Alexander and Davis & 2 & 2 \\
Nater, ICSNL 59 & 2 & 0 \\
Lyon, inchoatives & 2 & 0 \\
Matthewson and Redan & 1 & 1 \\
Wolfe & 1 & 0 \\
\bottomrule
\end{longtable}
\normalsize

\texttt{Kʷ əɬtəzétkʷu} reads as \texttt{Kʷəɬtəzétkʷu} at a score of 478.0, and that is Kʷəɬtəzétkʷu's own name broken in half in a paper that passes every check. \texttt{N ɬeʔkepmxcín} reads as \texttt{Nɬeʔkepmxcín} at 302.8 and \texttt{s ƛ̓eʔxímn.} as \texttt{sƛ̓eʔxímn.} at 244.1. A mark drawn at random from that paper's inventory lands on an attested form 0.000 of the time over 6600 draws.

The blanket repair cannot reach any of them and is right not to. The break follows \texttt{ʷ}, which is outside the joining set because it is a spacing modifier letter on the baseline and a space after one is normally a real boundary: including it welds thirteen pairs of real words in one paper. A blanket rule has to choose once for the whole corpus. This decides per site, from what the paper itself writes. It can take \texttt{Kʷ əɬtəzétkʷu} and leave \texttt{bəlkʷ `return'} alone.

The score separates them. Garcia's \texttt{w ɬ} scores $-1.4$ and is declined, because the runs \texttt{wɬ} would make sit below where a flat distribution puts them: the paper does not write words shaped that way. Nine of Garcia's 33 sites are declined on that test.

\textbf{These are candidate corrections, and none is applied.} Applying them rewrites the corpora and moves what every downstream measurement reads, which is a decision to take deliberately and measure. The table establishes that the damage is there, that it is larger than the checks can see, and that it is readable without a person opening a page.

\section{What the gold standard corpus measures}

The gold standard corpus that the language profiles are built from is 2314 lines and 16898 tokens across the eleven papers that have both an extractor and a profile to count under. 749 of those lines and 6762 of those tokens, 40 percent of it, come from the two papers set in the font that renumbers its codes.

Robertson's paper writes gold standard lines of its own and is not in that total: it is Thompson and Secwepemctsín written in Chinuk pipa, and there is no profile for it. Wolfe holds eighteen languages and writes no single-language file at all. The two Nater word lists, the Lyon inchoatives and the Kim reduplications each write one, and none of the four joins a profile, because a word list and a narrative in one language sit about as far apart as two languages do: Nater's voiceless list measures 0.7783 from the Nuxalk narrative corpus and Lyon's inchoatives 0.6436 from Nsyilxcən's, against between-language distances of 0.51 to 0.83. Joining one to a profile would move the profile toward the shape of a dictionary page.

The extractors for the two damaged papers read the drafted page text. Reading the extraction instead would put the font's alphabet in the corpus, with only the substitutions the font repair covers applied to it. Lindley and Lyon, whose table is transcribed from the pages in full and takes no repair, measures the difference token for token:

\tiny
\setlength{\tabcolsep}{2pt}
\begin{longtable}{@{}>{\raggedright\arraybackslash}p{\dimexpr\textwidth/3-2\tabcolsep\relax}>{\raggedright\arraybackslash}p{\dimexpr\textwidth/3-2\tabcolsep\relax}>{\raggedright\arraybackslash}p{\dimexpr\textwidth/3-2\tabcolsep\relax}@{}}
\toprule
 & reading the extraction & reading the page text \\
\midrule
corpus tokens the table holds & 1231 of 1878, 66\% & 2204 of 2535, 87\% \\
marks standing in front of their letter & 1832 & 54 \\
labialized consonants written with a plain \texttt{w} & 1630 & 42 \\
\bottomrule
\end{longtable}
\normalsize

The last two rows are for both Lyon papers together. Read from the extraction the corpus holds \texttt{iks’ma’yɬtím} where the page holds \texttt{iksm̓ay̓ɬtím}, and \texttt{sqwlqwlstwíxwtət} where the page holds \texttt{sqʷlqʷlstwíxʷtət}. Read from the page text it holds what the page holds.

What is left of the 13 percent is the inserted space: a word arrives in pieces and the corpus keeps the pieces. \texttt{incítxʷ} is in it as \texttt{incítx w}. The space repair cannot reach these, because the vocabulary it joins with comes from the interlinear's form column, which in this paper is the parse and carries morpheme hyphens the running text does not.

Every repair in the page-text path returns its line unchanged, because applying the mapping to text that has already been through it destroys the text: \texttt{P} to \texttt{ʔ} turns \texttt{Papers} into \texttt{ʔapers} and the gloss label \texttt{APPL} into \texttt{AʔʔL}. The language test on that path asks about the orthography itself, since the page text holds none of the damaged characters by construction.

\section{What a gold standard corpus is for}

The purity of both sources is the constraint everything else answers to. The target is the Salishan corpus and the reference corpus is English, and a boundary drawn between two distributions says nothing about the languages if either side is holding something other than what it claims.

\textbf{A paper holding eighteen languages has no single gold standard stream.} Every other extractor writes the lines of one language, because every other paper is one speaker telling one story, or one language under analysis. Wolfe cites lexical suffixes from published dictionaries of Sliammon, Sechelt, Squamish, three kinds of Halkomelem, three kinds of Straits, Klallam, Lushootseed, Twana, Tillamook, Quinault and two Tsamosan languages. Pouring those into one file builds a corpus of no language at all, and it would be a corpus the screening is then measured against, where that mistake cannot be recovered from. Its extractor writes its gold standard lines keyed by language, and the paper joins no profile.

Two further things are held out. A reconstruction is a form the author worked out and nobody has ever said, the same distinction as between a transcription and a segmentation line, applied to a proto-form. And the predicted reflexes of Tables 6 and 8, which give for each language the form expected if the reconstruction was stressed, the form expected if it was not, and the form attested; the extraction does not keep those three columns apart reliably. Nothing is lost, because the attested forms of both tables are printed again in examples (1) and (6) where the columns are unambiguous.

The hand extractions can also be turned on themselves. Each is verified against its own paper. They should agree with each other about what these languages look like, and one can be measured against the rest taken together as the reference corpus. A row that sits oddly against everything else verified is one of two things: a slip in the transcription, or a real property of that paper nobody has noticed. Both are worth opening the paper for, and neither is visible from inside one table. At the language level this runs: each dialect is measured against the other five pooled, the question the screening is actually asked, and all six stand off the pool. Per row it does not run, and per row it would name a paper to open.

Once each dialect's corpus is pure, the dialects become each other's reference corpora. A boundary drawn between Nsyilxcən and Nɬeʔkepmxcín is a claim about where one ends and the other begins, and it is worth drawing because of what sits near it: forms that two dialects share and a third has lost, which is where onomatopoeia that has gone out of one of them would show.

\textbf{One such boundary is drawn and scored against an answer the algorithm did not supply.} Lushootseed is not one dialect, and Mellesmoen and Kye's stress paper \cite{src:Mellesmoen-and-Kye-2026} labels every form it cites northern or southern. The hand extraction carries that label in the speaker column, a published fact from one paper, in one orthography, by one transcriber. The border test loads those labels, sets them aside, and only compares at the end. Comparing whole distributions fails at this size and says so: the two varieties hold 1469 and 2768 bytes against a method that resolves at 6707, and a blind partition scores exactly the majority class. Holding the concept fixed with the word web's gloss edge, flattening the pooled run counts to maximum entropy, and testing one run at a time does find it. The term is the stressed schwa \texttt{ə́}, southern, 45 against 6 over the 58 concepts both varieties name. Putting the border back at random 200 times, 1 of those random borders finds as much. The chapter \emph{How wrong it could be} carries the tables.

\textbf{That comparison has to run on a sound representation and not on the characters.} These orthographies write the same sound differently and the difference is arbitrary. The lateral fricative is \texttt{ɬ} in one paper here and \texttt{ł} in another. The glottal stop is \texttt{ʔ} in most and the digit \texttt{7} in the van Eijk papers. Labialization is a raised \texttt{ʷ} in some and a plain \texttt{w} in others. A model counting character bigrams reads two dialects that share a word as maximally far apart, because it is measuring the transcriber and not the speaker. Encoding each segment as its features and comparing those makes the distance a fact about the languages.

The plan runs in three stages. First the gold standard text corpus, which fixes what each dialect writes and which the hand extraction workflow exists to produce. Then a sound representation measured off recordings, where a segment is a set of binary features and the text space maps onto those features one written symbol at a time. Then a probability distribution over that representation, the thing an interdialect boundary is drawn on, with the word webs built on the encoded forms so that two dialects writing one word two ways arrive at one node.

The recordings the sound representation is measured from are three, all nɬeʔkepmxcín, all Bev Phillips. That is enough to build an encoding against and short of what a second dialect would take.

The word web is built on the written forms. It joins every form in the hand extractions to the forms it is related to, by three kinds of edge. A \textbf{concept} edge joins two forms whose glosses share a content word; it crosses an orthography, because the gloss is the only part of a form two papers wrote the same way. A \textbf{shape} edge joins two forms sharing a leading or trailing run of four characters, which in these languages is usually a shared root or affix. A \textbf{context} edge joins two forms written in the same section of one paper by one speaker. The web stands at 7062 forms and 373575 edges across nineteen groups. A dictionary multiplies the edges of a web keyed on glosses: 1275 entries each carrying an English gloss give the concept edge far more to join than the same number of words in running narrative.

The join between the web and the sound representation is not built: nothing takes a character out of the text space and hands back the features it stands for. The web's nodes are written forms, and two dialects writing one word two ways are two nodes. The text space is that join's other side, and it stays one definition for that reason.

The screening's calibration reports the separation as healthy against that corpus. Under one profile, 2285 known gold standard lines sit at a median of 15.15 against 1899 known English spans at 8.49, and a cut at 9.28 keeps 99.0 percent of the gold standard corpus while admitting 15.3 percent of the English. Under two profiles, 96 percent of the gold standard lines come back as language and 98 percent of the English spans as English. Those figures say the two profiles are far apart. They do not say the language profile is the language; the chapter \emph{The corpus under the instrument} gives what the six per-language profiles do, where they fail to separate, and why. The same measurement on the other papers is not comparable, because each goes through its own repair before the comparison, and several of those orthographies write labialization with a plain \texttt{w}.

\section{The binary sound representation}

The representation turns a recording into two bit fields per 10 ms frame. Four facts about hearing set its shape.

The waveform on its own is not enough. A vowel is a harmonic series under an envelope, and two recordings of one vowel share the envelope while the harmonics sit wherever the speaker's pitch put them. The analysis is spectral and the envelope is smoothed off the source before anything is compared.

What a person hears is not what the microphone recorded, and it differs between listeners by physiology and by how much hearing they have lost. The weighting takes an audiogram, the six frequencies a hearing test is run at with the loss in dB at each; the loss comes off a band before loudness, and a band a listener cannot hear stops reaching the code. The default is a listener who has lost none.

Some sounds are the same sound at different frequencies. A phone said by a small speaker and by a large one differs by a scaling of the whole spectrum. The 64 bands are log spaced, and that scaling is then a shift along the band axis, and the magnitude of a Fourier transform along that axis is unchanged by the shift. That transform is the segment field.

And some sounds that differ only in frequency mean different things. Stress carries an acute in every orthography in this corpus, and a field made shift invariant would put a stressed form and an unstressed one at one address. The prosody field is separate for that reason and holds loudness, pitch against the speaker's own median, how that pitch is moving, and voicing.

\textbf{Bringing the sample to maximum entropy makes the delta readable.} Each bit is cut at its own column's median over the recording and splits the frames in half, and the segment columns are rotated onto their principal axes first so that no bit restates another. All 24 segment bits come out set on 0.500 of frames on all three recordings, and all 64 of the six bit states and all 256 of the eight bit states are reached on all three. A recording that used those states evenly would be noise, and the distance from that flat reference is the structure in it.

\tiny
\setlength{\tabcolsep}{2pt}
\begin{longtable}{@{}>{\raggedright\arraybackslash}p{\dimexpr\textwidth/6-2\tabcolsep\relax}>{\raggedright\arraybackslash}p{\dimexpr\textwidth/6-2\tabcolsep\relax}>{\raggedright\arraybackslash}p{\dimexpr\textwidth/6-2\tabcolsep\relax}>{\raggedright\arraybackslash}p{\dimexpr\textwidth/6-2\tabcolsep\relax}>{\raggedright\arraybackslash}p{\dimexpr\textwidth/6-2\tabcolsep\relax}>{\raggedright\arraybackslash}p{\dimexpr\textwidth/6-2\tabcolsep\relax}@{}}
\toprule
Recording & Frames & 8 bits & 10 bits & 12 bits & Prosody \\
\midrule
Givens and Hall, ICSNL 58 & 23489 & 0.2157 & 0.2538 & 0.3161 & 0.3708 \\
Hall and Phillips, ICSNL 59 & 83249 & 0.2364 & 0.2863 & 0.3207 & 0.3498 \\
Hall and Phillips, ICSNL 60 & 134997 & 0.1859 & 0.2200 & 0.2519 & 0.4594 \\
\bottomrule
\end{longtable}
\normalsize

Each number is the total variation between the codes that width of field took and the flat distribution over its states. The width is part of the figure because the field has to be sampled to be read. At 12 bits the longest recording puts 33 frames in each of 4096 states; at 24 bits it puts 133820 codes in 16.8 million states, which measures the frame count and not the recording.

All three recordings came from the nɬeʔkepmxcín Lab and all three are Bev Phillips. These numbers are the method working and not a comparison of dialects. A second dialect's recording turns them into a reading.

\section{Audio sources}

\tiny
\setlength{\tabcolsep}{2pt}
\begin{longtable}{@{}>{\raggedright\arraybackslash}p{\dimexpr\textwidth*2/3-2\tabcolsep\relax}>{\raggedright\arraybackslash}p{\dimexpr\textwidth/3-2\tabcolsep\relax}@{}}
\toprule
Recording & Size \\
\midrule
Hall and Phillips, xʷíʔ kʷ páq (You Will Be Sorry), ICSNL 59 \cite{src:Hall-and-Phillips-2024} & 13.4 MB \\
Givens and Hall, The Moon and the Birchbark Canoe, ICSNL 58 \cite{src:Givens-and-Hall-2023} & 7.5 MB \\
Hall and Phillips, ɬ cutés us ɬ qəɬmín ɬ tmíxʷ, ICSNL 60 \cite{src:Hall-and-Phillips-2025} & 17.1 MB \\
\bottomrule
\end{longtable}
\normalsize

The third is Bev Phillips reading the story whose paper has an extractor, and it is a test oracle for that extraction as well as a source.

Recordings named in papers and not published with them:

\begin{itemize}
\item Qwa7yán'ak, Graveyard Valley narrative. Just over half an hour, recorded by Henry Davis at Nxwísten on 7 July 2025.
\item K̓weswapáw̓ (Linda Redan). Three minutes twenty-eight seconds, told over Zoom on 31 October 2025. Audio and video are held by her.
\item George Lezard, 1966, recorded by Randy Bouchard.
\item Nellie Guitterez, 1978 or 1979, recorded by Yvonne Hébert at an Elders' Gathering in Quilchena.
\end{itemize}

\section{Unread papers}

152 extracted paper texts are held, and 146 of the PDFs they came from. Eighteen have an extractor and a hand extraction, and seven more have a hand extraction alone. Six of the 146 are in the font class above, and their texts are the font's encoding. Fifty of the papers are about Lushootseed or Skagit.

The screening reads each paper's front matter for the names the literature uses for each language and attributes a paper where it names one and no other. Over 154 texts it attributes 122 and leaves 32, which are the comparative papers and are correctly left. Lushootseed leads at 43 papers, Nuxalk follows at 32.

Counting how much of a language each unread paper actually holds gives a different order, and it is the order to read by. Lushootseed has 40 unread papers and 593 distinct language tokens across all of them together, which is less than one Nater appendix. Most of the Lushootseed literature is older and writes the glottal stop as \texttt{?} and the schwa as \texttt{a} or \texttt{E}, and a mark set built on the modern orthography reads those papers as holding no language at all. Comox has 8 unread papers and 1899 tokens, St'át'imcets 10 and 1297, Nuxalk 29 and 1046.

The count of papers per language is not the count of language per language. What Lushootseed needs is not another paper but the treatment Hilbert's 1983 paper and Hilbert and Hess already have: a per-paper mark set for the orthography that paper is actually written in.

\section{Verification sources}

Some sources test a transformation instead of supplying text. Each answers a question the code that made the change cannot answer about itself.

\begin{itemize}
\item Lyon's inchoatives paper \cite{src:Lyon-2025}, a later paper on Nsyilxcən whose extraction kept its characters. It tests the font substitution table, which moves attested tokens from 1 of 3599 to 811 on one paper and from 2 of 4332 to 965 on the other, and the word list recovered from the interlinear, where none of the broken forms has its pieces attested apart and 4 percent and 36 percent are attested once joined.
\item Each paper's own second printing. Hall and Phillips, Garcia, and both Lyon papers print their stories twice, once as running text and once word by word. One printing checks a reading of the other.
\end{itemize}

\section{Community sources}

\begin{itemize}
\item Puyallup Tribal Language Program, \texttt{https://www.puyalluptriballanguage.org/}
\item The Salish Institute, \texttt{https://www.thesalishinstitute.com/salish-language}
\end{itemize}

\section{Archives}

\begin{itemize}
\item ICSNL volumes, 1967 to present, every paper on one page: \texttt{https://lingpapers.sites.olt.ubc.ca/icsnl-volumes/}.
\item UBCWPL, volumes 33 (1998) to 61 (2026): \texttt{https://lingpapers.sites.olt.ubc.ca/}.
\item ICSNL archives, volumes 1 to 32, the Kinkade Collection: \texttt{https://blogs.ubc.ca/icsnl/icsnl-archives-and-precedings/}.
\item nɬeʔkepmxcín Lab, the source of the three recordings: \texttt{https://blogs.ubc.ca/nlab/projects-publications/}.
\item ICSNL 50 as one volume: \texttt{https://lingpapers.sites.olt.ubc.ca/files/2018/01/ICSNL2015-fullonline.pdf}.
\end{itemize}
